\section{化学方程式} 

化学方程式可以直接采用数学式输入，例如: 
三硝基甲苯(TNT)~C$_6$H$_2$CH$_3$(NO$_2$)$_3$
为白色或苋色淡黄色针状结晶，无臭，有吸湿性。

很明显，将化学方程式作为数学公式输入很复杂，且十分笨拙。所以本模版引入了mhchem宏包将问题简化。~\verb|\ce|~命令用来输入化学方程式。 如：醋中主要是 \ce{H2O}，含有 \ce{CH3C00-}。\ce{^{277}_{90}Th}元素具有强放射性。

 化学反应式例子如下：
 \begin{gather} 
   \ce{2H2 + O2 ->[\text{燃烧}] 2H2O} \\
   \ce{N2 + 3H2 <=>T[高温、加压][催化剂] 2NH3}
 \end{gather}

有机化学式的书写，先简单介绍一下~chemfig~方向的定义，如图~\ref{fig:direction}：

\begin{figure}[ht]
  \centering
    \includegraphics[width=0.2\textwidth]{direction.png}
    \caption{chemfig方向的定义}
    \label{fig:direction}
\end{figure}

举个乙烷的例子：

\begin{center}
  \chemfig{C(-[2]H)(-[4]H)(-[6]H)-C(-[2]H)(-[6]H)-H}
\end{center}

上述乙烷代码如下：
{\verb|\chemfig{C(-[2]H)(-[4]H)(-[6]H)-C(-[2]H)(-[6]H)-H}|}
代码中 (-[X]Y) 内的数字 X 就表示了内容 Y 的位置，其中中括号（[ ]）的位置通常紧跟在化学键（-、= 等）之后。例如2表示的就是向上的方向，4 表示向左，6 表示向右。

再画一个苯环:
\begin{center}
  \chemfig{*6(-=-=-=)}
\end{center}